(1) 中心引力加速度 \(g(r)=GM/r^2\)，近端相对卫星中心的差 \(|g(r-R)-g(r)|\simeq2GMR/r^3\)。注意两端之间的全差是 \(4GMR/r^3\)，本题判据始终采用单侧半径尺度。
(2) 卫星质量 \(m_s=4\pi\rho R^3/3\)，表面自引力 \(g_s=4\pi G\rho R/3\)。忽略内压、旋转和材料强度，以 \(2GMR/r^3\gtrsim4\pi G\rho R/3\) 估计破裂；这是数量级判据，不是自由表面精确平衡。
(3) 等号得 \(r_c=[3M/(2\pi\rho)]^{1/3}\)；若 \(M=4\pi\rho_pR_p^3/3\)，则 \(r_c=R_p(2\rho_p/\rho)^{1/3}\)。卫星半径 \(R\) 消去，量纲为长度，显示自引力和潮汐均正比 \(R\)。
(4) 圆轨道 \(\Omega^2=GM/r^3\)，临界时 \(\Omega_c^2=2\pi G\rho/3\)。定义 \(t_{\rm dyn}=\sqrt{3/(4\pi G\rho)}\)，则 \(\Omega_ct_{\rm dyn}=1/\sqrt2\)。真实洛希系数受同步自转、液体变形与压强分布影响，本式不得冒充精确常数。
